\input{formattingHeader}

\titleandheader{Checking the OpenLCB Stream Transport Protocol}

\begin{document}
\maketitle
\thispagestyle{firststyle}

\introductionCaveats
    {https://nbviewer.org/github/openlcb/documents/blob/master/standards/StreamTransportS.pdf}
    {Stream Transport Standard}.

\section{Stream Transport Procedure}

\checkProcedure{Stream Checking}

A node which does not self-identify in PIP that it supports
Stream Transport will be considered to have passed these checks.
\pipsetFootnote

The checker acts as the Source node (initiating streams to the DBC).
The DBC is always the Destination node.  For nodes that do not support
streams, we verify proper rejection.

This plan does not check:
\begin{enumerate}
\item Streams initiated by the DBC as source.  Config-mem stream checks
    exercise that path separately.
\item CAN multi-frame assembly of 12-byte Initiate Requests.  The st040 check
    sends a 12-byte request but does not explicitly verify CAN framing.
\item Gateway buffer reduction (not checkable from an end node).
\item Early Proceed (implementation optimization, not a conformance requirement).
\end{enumerate}

All checks validated against PR \#189 (openlcb/documents), the authoritative
spec revision.

\subsection{Stream Initiate Request/Reply}

This checks Standard sections 4.1, 4.2 and 7.1.

The checker sends a Stream Initiate Request with a proposed buffer size of 256
and Source Stream ID 0x01.  It then checks:
\begin{enumerate}
\item That a Stream Initiate Reply or Optional Interaction Rejected is received.
\item If a Reply: length is at least 6 bytes, SID is echoed, DID is not 0xFF,
    flags and buffer size are self-consistent.
\item If accepted, the stream is closed with Stream Data Complete.
\item A 30-second stability soak monitors for unexpected reboots.
\end{enumerate}

\subsection{Reject When Unsupported}

This checks Standard section 7.1.

If the DBC does not advertise Stream Transport in PIP, the checker sends a
Stream Initiate Request and verifies the DBC rejects it with OIR or an error
reply.

\subsection{Buffer Negotiation}

This checks Standard section 4.2.

The checker opens streams with proposed buffer sizes of 64 and 1024.  For each:
\begin{enumerate}
\item Verifies the reply buffer size does not exceed the proposed size.
\item Verifies the reply buffer size is greater than zero if accepted.
\end{enumerate}

\subsection{Minimum Buffer Size}

This checks Standard section 4.2.

The checker sends a Stream Initiate Request with buffer size 1.  Verifies
acceptance with size 1 or a valid rejection.

\subsection{Maximum Buffer Size}

This checks Standard section 4.2.

The checker sends a Stream Initiate Request with buffer size 0xFFFF.  Verifies
the reply buffer does not exceed 0xFFFF and the node negotiates down to its
capability.

\subsection{Unknown Content UID}

This checks Standard section 6.2.

The checker sends a 12-byte Stream Initiate Request with a fabricated Content
UID.  Accepts rejection with 0x1040 (unimplemented) or acceptance (both
spec-legal).

\subsection{Initiate Reply Format}

This checks Standard sections 4.1 and 4.2.

The checker sends an initiate with SID 0x42 and validates every field in the
reply:
\begin{enumerate}
\item Length at least 6 bytes.
\item Byte 4 echoes SID 0x42 exactly.
\item Byte 5 (DID) in range 0x00--0xFE.
\item If accept: flags 0x0000 or 0x8000, buffer greater than 0 and not
    exceeding proposed.
\item If reject: error code 0x1xxx or 0x2xxx, buffer 0.
\end{enumerate}

\subsection{Data Send and Proceed}

This checks Standard sections 4.3, 4.4 and 6.3.

The checker opens a stream, sends exactly MaxBufSize bytes, and verifies a
Stream Data Proceed is received with correct SID and DID.  The stream is then
closed.

\subsection{Data Proceed Flow Control}

This checks Standard section 6.3.

The checker opens a stream and performs two consecutive buffer cycles:
\begin{enumerate}
\item Send MaxBufSize bytes, wait for Proceed.
\item Send another MaxBufSize bytes, wait for second Proceed.
\end{enumerate}
Verifies no unsolicited errors during transfer.

\subsection{Stream Completion}

This checks Standard section 6.4.

The checker opens a stream, sends a partial buffer (less than MaxBufSize),
sends Stream Data Complete, verifies no errors, then reopens a new stream to
confirm resources were freed.

\subsection{Terminate Due To Error}

This checks Standard section 6.4.

The checker opens a stream, sends a Terminate Due To Error, verifies the node
remains stable, then reopens a stream to confirm recovery.

\subsection{Post-Stream Stability}

This checks general robustness after stream operations.

The checker opens and closes a stream, then monitors the node for 30 seconds,
pinging every 10 seconds with Verify Node ID Addressed.  Watches for
unexpected Initialization Complete messages.

\subsection{OIR Fallback}

This checks Standard section 7.1.

When Stream Transport is in PIP, the checker sends a Stream Initiate Request
and accepts either a Stream Initiate Reply or OIR as valid responses.

\subsection{Concurrent Streams, Same Source Node}

This checks Technical Note section 1 and section 2.3.1.

The checker opens two concurrent streams (SID 0xA0 and 0xB0) to the DBC.
If the second is rejected, the check passes with info (single-stream node).
Otherwise:
\begin{enumerate}
\item Sends interleaved data on both streams.
\item Verifies each Proceed carries the correct SID/DID for its stream.
\item Repeats for a second cycle.
\item Closes both streams.
\end{enumerate}

\subsection{Concurrent Streams, Different Source Nodes}

This checks Technical Note section 2.3.1.

The checker claims a second CAN alias (spoofed source node) and opens
concurrent streams from both the real and spoofed nodes to the DBC, using the
same SID (spec-legal since different source nodes).  Verifies the DBC tracks
streams by (Source Node ID, SID) and does not confuse traffic from different
sources.

\subsection{Sustained Transfer}

This checks Standard section 6.3.

The checker opens a stream and runs 20 consecutive send/proceed cycles with
zero delay between them.  Verifies the DBC does not leak memory, lose window
count, or stall.  Reopens a stream afterward to confirm recovery.

\subsection{Minimum Chunk Send}

This checks Standard section 4.3.

The checker sends MaxBufSize bytes as individual 1-byte payload Stream Data
Send messages, maximizing the number of CAN frames per buffer window.
Stresses the receive path and any internal reassembly.

\subsection{Concurrent Sustained}

This checks Standard section 6.3 and Technical Note section 1.

Two concurrent streams both doing 10 sustained multi-cycle transfers
simultaneously.  Stresses total buffer pool allocation.

\subsection{Asymmetric Buffers}

This checks Standard section 4.2.

Two concurrent streams with different proposed buffer sizes (64 and 1024).
Verifies the DBC tracks different window sizes and sends Proceed at the
correct boundary for each stream independently.

\subsection{Negotiation Enforcement}

This checks Standard section 4.2.

The checker proposes a large buffer (1024).  The DBC may negotiate down.
The check sends exactly the negotiated amount and confirms Proceed arrives at
the right boundary, not at the proposed size.

\subsection{Zero-Payload Flush}

This checks Standard section 6.3 and Technical Note section 2.4.3.

The checker sends half a buffer window, then a zero-payload Stream Data Send
(just the DID byte), then the remaining half.  Verifies Proceed arrives after
the full MaxBufSize payload bytes (the flush does not count toward the buffer
window).  Repeated for 2 cycles.

\subsection{Interleaved Data Send}

This checks Standard section 7.2 and Technical Note section 2.3.1.

Two concurrent streams receive interleaved Stream Data Send messages within
the same buffer window.  The checker alternates 32-byte chunks (A, B, A, B)
until both windows are full.  Forces the DBC to track each stream's byte count
independently.  Repeated for 2 cycles.

\subsection{Suggested DID}

This checks Standard section 4.1.

The checker sends a Stream Initiate Request with a suggested DID of 0x42
(byte 0 of the request).  Verifies the reply DID is in range 0x00--0xFE.
The DBC may honor the suggestion or assign its own; both are valid.
Then reopens with DID 0xFF (no suggestion) to confirm both paths work.

\subsection{Complete with Byte Count}

This checks Standard section 6.4.

This checks both forms of Stream Data Complete:
\begin{enumerate}
\item With the optional 4-byte total byte count appended.
\item Without the byte count (2 bytes only).
\end{enumerate}
For each form, sends one full buffer, waits for Proceed, sends Complete,
verifies no errors, and reopens a stream to confirm resources were freed.

\subsection{Reject Then Retry}

This checks Standard section 6.2.

The checker opens streams until one is rejected (up to 4 attempts).  If no
rejection occurs, the check passes (deep pool node).  Otherwise:
\begin{enumerate}
\item Closes one open stream to free resources.
\item Retries the rejected initiate.
\item Verifies the retry succeeds.
\item Closes all streams.
\end{enumerate}

\section{Compliance Matrix}

The following table maps spec requirements to the checks that cover them.

\subsection{Covered Requirements}

\begin{tabular}{|p{8cm}|p{5cm}|}
\hline
\textbf{Spec Requirement} & \textbf{Check(s)} \\
\hline
S4.1 Stream Initiate Request format & st010, st050, st040 \\
\hline
S4.1 Suggested DID in Initiate Request & st190 \\
\hline
S4.2 Stream Initiate Reply format, error codes, SID echo, DID range & st010, st050, st060 \\
\hline
S4.2 Accept: buf $\leq$ proposed, buf $>$ 0 & st030, st035, st037, st155 \\
\hline
S4.2 Reject: permanent/temporary error codes, buf = 0 & st020, st050 \\
\hline
S4.3 Stream Data Send with DID byte 0 & st060, st065, st070 \\
\hline
S4.4 Stream Data Proceed: SID/DID correct & st060, st065, st100, st140 \\
\hline
S4.6 Stream Data Complete closes stream & st070, st080 \\
\hline
S5 States: Initiated $\rightarrow$ Open transition & st010 \\
\hline
S5 States: Open $\rightarrow$ Closed via Complete & st070 \\
\hline
S5 States: Open $\rightarrow$ Closed via TDE & st080 \\
\hline
S6.1 Opening with announcement (via config-mem) & (memory config stream checks) \\
\hline
S6.2 Opening without announcement (Content UID) & st040 \\
\hline
S6.2 Unknown Content UID $\rightarrow$ 0x1040 or accept & st040 \\
\hline
S6.2 Reject then retry after freeing resources & st210 \\
\hline
S6.3 Data transfer: send up to MaxBufSize, wait for Proceed & st060, st065 \\
\hline
S6.3 Proceed extends window by MaxBufSize & st065, st120, st155 \\
\hline
S6.4 Stream Data Complete closes stream & st070 \\
\hline
S6.4 Stream Data Complete with optional byte count & st200 \\
\hline
S6.4 TDE closes stream & st080 \\
\hline
S7.1 CAN Frame Type 7: DID in byte 0, payload in bytes 1--7 & st060, st110, st125 \\
\hline
S7.2 Message-level interleaving across concurrent streams & st180 \\
\hline
TN1 Up to 255 concurrent streams per node pair & st100 (2 streams) \\
\hline
TN2.3.1 Stream identified by (SrcNodeID, DstNodeID, SID) & st110 \\
\hline
TN2.4.3 Zero-payload flush & st170 \\
\hline
Post-operation stability & st085 \\
\hline
OIR as valid fallback & st090 \\
\hline
Resource recovery after close & st070, st120, st140 \\
\hline
\end{tabular}

\subsection{Not Covered}

\begin{tabular}{|p{4cm}|p{2.5cm}|p{6cm}|}
\hline
\textbf{Gap} & \textbf{Spec Ref} & \textbf{Notes} \\
\hline
12-byte Initiate Request CAN multi-frame & S7.2 & st040 sends it but does not explicitly verify CAN framing. \\
\hline
Gateway buffer reduction & S6.5 & Not checkable from an end node. \\
\hline
Proceed sent early & TN2.4.4, S8.4 & Implementation optimization, not a conformance requirement. \\
\hline
Stream from DBC as source & S6.1, S6.3 & Config-mem stream checks cover this path. \\
\hline
\end{tabular}

\end{document}
